\documentclass[12pt]{article}
\usepackage[utf8]{inputenc}
\usepackage[russian]{babel}
\usepackage{amsmath}
\usepackage{amssymb}
\usepackage{geometry}
\geometry{a4paper, margin=2cm}

\begin{document}

\section*{Задача 3.1.8}

\textbf{Условие:} Проверить, является ли входное число $n\geq2$ простым. Вывести $1$, если число простое, и $0$ в противном случае.

in: $n$\\
out: $r\in\{0,1\}$

\section*{Псевдокод}

\begin{verbatim}
SMALL_PRIMES = (
    2, 3, 5, 7, 11, 13,
    17, 19, 23, 29, 31, 37
)
MR_BASES = (
    2, 325, 9375, 28178,
    450775, 9780504, 1795265022
)
LIMIT = 1 << 64

def mod_pow(a_1, d_1, n):
    result_1 = 1

    while d_1 > 0:
        if d_1 % 2 == 1:
            result_1 = (result_1 * a_1) % n

        a_1 = (a_1 * a_1) % n
        d_1 //= 2

    return result_1

def miller_rabin(n):
    d = n - 1
    s = 0
    while d % 2 == 0:
        d //= 2
        s += 1
    for a in MR_BASES:
        a %= n
        if a == 0:
            continue
        res = mod_pow(a, d, n)
        if res == 1 or res == n - 1:
            continue
        for _ in range(s - 1):
            res = res * res % n
            if res == n - 1:
                break
        else:
            return False
    return True

def is_prime(n):
    if n < 2:
        return False
    for p in SMALL_PRIMES:
        if n % p == 0:
            return n == p
    res = miller_rabin(n)
    if n < LIMIT:
        return res
    if not res:
        return False
    d = 41
    while d * d <= n:
        if n % d == 0 or n % (d + 2) == 0:
            return False
        d += 6
    return True
\end{verbatim}

\section*{Сложность}

Рассмотрим 2 случая:
\begin{enumerate}
    \item $n < 2^{64}$. Тогда тест Миллера-Рабина работает точно. Худшим случаем будет, когда не будет пройдена проверка на малых простых числах.
    \item $n \ge 2^{64}$. Тогда тест работает неточно. Худшим случаем будет, когда проверка не отсеится на малых простых числах и если тест Миллера-Рабина даст положительный ответ. Тогда ответ нужно будет искать с помощью алгоритма 6k+1.
\end{enumerate}

\begin{enumerate}
    \item is\_prime без miller\_rabin
    \begin{itemize}
        \item Проверка входа и инициализация таблиц - $46$ шагов
        \item Цикл FOR\_SMALL: если ни одно из $12$ малых простых чисел не делит $n$, каждая итерация занимает $15$ шагов. Выход из цикла - $3$ шага. Итого: $12\cdot15+3=183$ шага
        \item После возврата из miller\_rabin проверка границы $2^{64}$ занимает $3$ шага. Если $n<2^{64}$, обработка результата и вывод занимают ещё $4$ шага при ответе $0$ или $5$ шагов при ответе $1$
        \item Если $n\geq2^{64}$ и тест Миллера-Рабина вернул $0$, проверка результата, вывод и остановка занимают $4$ шага после проверки границы. Если тест вернул $1$, запуск дополнительного перебора - $8$ шагов
        \item Пусть $q$ - число итераций дополнительного цикла, значит $d=41+6j$ и $d^2\leq n$. Тогда:
        \[
            q=\left\lfloor\frac{\lfloor\sqrt n\rfloor-41}{6}\right\rfloor+1
        \]
        Каждая итерация без найденного делителя занимает $23$ шага, проверка выхода - $4$ шага, вывод и остановка - $2$ шага. Поэтому цикл - $23q+6$ шагов. Ну и тогда вся часть после положительного результата Миллера-Рабина - $23q+14$ шагов.
    \end{itemize}

    \item miller\_rabin без mod\_pow
    \begin{itemize}
        \item Инициализация d = n - 1, s = 0.
        В цикле while d \% 2 == 0 значение $d$ делится на $2$,
        а $s$ увеличивается на $1$. Цикл выполняется $s$ раз, после
        чего $d$ становится нечётным. Эта часть - $14+13s$ шагов.
        \[
            s(n)=\max\{k\geq0:(n-1)\bmod 2^k=0\}
        \]
        Сначала $d = n - 1$, значит $s\le \lfloor\log_2(n-1)\rfloor$.

        \item Для $n\geq2^{64}$: $\forall i\in MR\_BASES \; i\not\equiv_n0$.
        Для каждого MR\_BASES проверка цикла, получение MR\_BASES и
        вычисление $i \bmod n$ - $15$ шагов. Для семи $i\in MR\_BASES$ -
        $105$ шагов

        \item Пусть $r_i:=a_i^d\bmod n$. После возврата из
        mod\_pow, проверки - $P_i$ шагов:
        \[
            P_i=\begin{cases}
                4, & r_i=1\\
                8, & r_i=n-1\\
                17t_i+7, & r_i\notin\{1,n-1\}
            \end{cases}
        \]
        В третьем случае $t_i$ - число $res * res \% n$,
        после которых $res == n-1$. По программе $t_i<s-1$

        \item Переход к следующему MR\_BASES - $4$ шага, значит
        $28$ шагов для всех MR\_BASES. Финальная проверка цикла
        и сохранение результата - $5$ шагов

        \item Таким образом, для выполнения всех семи 
        MR\_BASES без учёта mod\_pow:
        \[
            T_{MR\setminus mod\_pow}
            =14+13s+105+\sum_{i=1}^{7}P_i+28+5
            =152+13s+\sum_{i=1}^{7}P_i
        \]
    \end{itemize}

    \item mod\_pow
    \begin{itemize}
        \item Пусть $h=\lfloor\log_2 d\rfloor+1=|d|$, $w$ - число единиц в $d$
        \item Инициализация - $2$ шага, выходная проверка цикла - $3$ шага
        \item Внутренний цикл выполняется $h$ раз
        \begin{enumerate}
            \item Для нулевого бита $d$ - $24$ шага
            \item Для единичного бита - $31$ шаг
        \end{enumerate}
        \item Поэтому один вызов mod\_pow:
        \[
            T_{mod\_pow}(n)
            =2+24(h-w)+31w+3
            =24h+7w+5
        \]
        Тогда для всех MR\_BASES - $168h+49w+35$ шагов.
    \end{itemize}
\end{enumerate}

В худшем случае предварительные проверки - $46+183=229$
шагов. Для $n<2^{64}$ максимальное число шагов:
\[
T_1(n)=424+13s+168h+49w+\sum_{i=1}^{7}P_i
\]
Как показано выше $s=O(\log n)$, $h=O(\log n)$, $w\leq h$ и
$t_i\leq s-1$, значит $T_1(n)=O(\log n)$.

Для $n\geq2^{64}$ худший случай:
\[
T_2(n)=433+13s+168h+49w+\sum_{i=1}^{7}P_i+23q
\]
где
\[
q=\left\lfloor
\frac{\lfloor\sqrt n\rfloor-41}{6}
\right\rfloor+1=O(\sqrt n)
\]
Остальные слагаемые $O(\log n)$, поэтому
сложность во втором случае $T_2(n)=O(\sqrt n)$.

\end{document}